\section{Conclusion}
\label{sec:conclusion}

We presented EWL, a progress-adaptive sample weighting method for LoRA
fine-tuning that tracks per-sample loss velocity via EMAs and gates the
resulting signal with a LoRA gradient proxy, requiring no meta-gradients,
no held-out data, and no inference cost.
Two findings stand out.
First, the gradient proxy is essential: ablating it causes a $-6.76\%$ drop on
FGVC-Aircraft, as ungated reweighting destabilises training on narrow decision
boundaries.
Second, EWL's noise robustness is dataset-conditional: it improves over SFT
by $+4.3$--$6.8\%$ under noise on Aircraft, where inter-class similarity
sustains a wide adaptation frontier, but is harmful ($-0.9\%$ to $-6.6\%$) on
CUB-200 and Stanford Dogs where the frontier is narrow.
On clean data EWL matches SFT within standard deviation.
EWL is therefore a principled but conditional method, best suited to
fine-grained, noise-prone tasks; its applicability is tied directly to
adaptation frontier width.
Future work includes vectorising the per-sample loop onto GPU tensors to cut
the $5.9\times$ overhead, combining velocity with class-aware reweighting,
and extending the frontier hypothesis to language models and other PEFT
methods.
